% !TEX encoding = UTF-8 Unicode
\usepackage[a4paper, top=25mm, bottom=25mm, left=25mm, right=25mm]{geometry}
\usepackage{microtype} % Tweak font spacing for aesthetics
%\usepackage{cmbright} % Font
\usepackage[utf8]{inputenc}
\usepackage[T1]{fontenc} 
\usepackage[italian]{babel} 
\usepackage[colorlinks=true, allcolors=black]{hyperref}
\usepackage{subfigure} % Multiple figures environment
\usepackage{amsmath,amssymb}
\usepackage{graphicx} % Include graphics
\usepackage{booktabs} % Professional tables
\usepackage[dvipsnames]{xcolor}  % Advanced colours
\usepackage{soul} % Highlight text and other text things
\usepackage{caption} % Manage captions in floats
\usepackage{textcomp} % Includes euro symbol
\usepackage{ifthen} % Use of conditional \if
\usepackage{fmtcount} % Convert integers to words
\usepackage[per-mode = symbol]{siunitx}  % To use SI units
\usepackage[version=4]{mhchem}  % For writing chemical symbols
\usepackage{totcount} % Display points counter at beginning 
\usepackage{xstring} % To manipulate strings
\usepackage{multicol} % Multiple column environment (FS)
%\usepackage{lipsum,turnthepage}
\usepackage{wasysym} % Symbols for multiple choice questions
\usepackage{fontawesome}
\usepackage{fontawesome}
\usepackage[most]{tcolorbox} % to circle the text of the email
\usepackage{multicol}

% DOCUMENT DEFINITIONS
\pointsinrightmargin  % Points on the right margin
\pointformat{\bfseries(\themarginpoints)}  % Points in bold and bracket style
\renewcommand{\partlabel}{\bfseries{(\thepartno)}} % Bold part numbers
\renewcommand{\subpartlabel}{\bfseries{\thesubpart)}} % Bold part numbers
\pagestyle{headandfoot} % Solution header and footer
\header{}{}{} % Solution header
\footer{}{}{\theCourse{} -- \theExamDate{} -- Gruppo \theGroupID{}} 
%\graphicspath{{img/}} % Directory where pictures are stored
\definecolor{myRed}{RGB}{184,0,1} % Define new colour
\definecolor{myGrey}{RGB}{211,211,211}  % Define new colour
\definecolor{verylightgray}{RGB}{241,241,241}  % Define new colour
\checkedchar{\large $\CIRCLE$}
\setlength{\fboxrule}{0pt}

% VERSION TRACKING
\newcommand{\theVersion}{1.1}
\newcommand{\theVersionDate}{13-Dec-2021}

% COURSE
\def\theExamEmail{esame.mirpalab@gmail.com}
\newcommand{\theCourseCode}{CODE12345}
\newcommand{\theDiet}{Diet} % EDIT - Exam diet (Main OR Resit)
\newcommand{\examOC}{Open}
\newcommand{\useCalculators}{Yes}
\newcommand{\formulaSheet}{Yes}
\newcommand{\hasSections}{Yes}
\newcommand{\separateBooks}{No} % Answer sections in separate books
\newcommand{\numberQuestions}{All}
\newcommand{\withSolutions}{Yes}

% NEW COMMANDS
\newtotcounter{pointsCounter} % Counter for total number of points in parts
\newtotcounter{numberOfSections} % Counter for number of sections
\newtotcounter{numberOfQuestions} % Counter for number of questions

% EXAM STRUCTURE - SECTIONS AND QUESTIONS
\ifthenelse{\equal{\hasSections}{Yes} \OR \equal{\hasSections}{Y} \OR \equal{\hasSections}{yes} \OR \equal{\hasSections}{y} \OR \equal{\hasSections}{YES}}
% \hasSections THEN clause
{\newcommand{\sectionStart}{
    \setcounter{question}{0} % (Re)initiate questions counter
    \setcounter{figure}{0} % Initiate figure counter
    \setcounter{table}{0} % Initiate table counter
    \qformat{\ifthenelse{\equal{\thequestion}{1}}
        {\fbox{\parbox{16cm}{\textbf{\large SEZIONE \Alph{numberOfSections} \hfill\vrule depth 0.5em width 0pt \\ 
        Quesito \Alph{numberOfSections}\thequestion}} \hfill\vrule depth 1.5em width 0pt}}
        {\fbox{\parbox{16cm}{\large \textbf{Quesito \Alph{numberOfSections}\thequestion}} \hfill\vrule depth 0.8em width 0pt}}
    }
    \stepcounter{numberOfSections} % Increment sections counter
}}
% \hasSections ELSE clause
{   \setcounter{question}{0} % (Re)initiate questions counter
    \setcounter{figure}{0} % Initiate figure counter
    \setcounter{table}{0} % Initiate table counter
    \qformat{
        {\fbox{\parbox{16cm}{\large \textbf{Quesito \thequestion}} \hfill\vrule depth 0.8em width 0pt}}
    }
}

\newcommand{\callQuestion}[2]{ % Call new question
    \cleardoublepage \input{#1} % Start question in new page
    \addtocounter{pointsCounter}{#2} % Add points to counter
    \addtocounter{numberOfQuestions}{1} % Increment questions counter
    \stepcounter{figure} % Increment figure counter
    \stepcounter{table} % Increment table counter
}

\newcommand{\callFormulaSheet}[1]{
    \ifthenelse{\equal{\formulaSheet}{Yes} \OR \equal{\formulaSheet}{Y} \OR \equal{\formulaSheet}{yes} \OR \equal{\formulaSheet}{y} \OR \equal{\formulaSheet}{YES}}{\cleardoublepage \input{#1}}{}
}

\newcommand{\headingFormulaSheet}{
    \pagestyle{empty}
    \setlength{\textwidth}{180mm}
    \setlength{\textheight}{265mm}
    \setlength{\topmargin}{-20mm}
    \setlength\oddsidemargin{-12mm}
    \setlength\evensidemargin{-12mm}
    \setlength{\columnsep}{10mm}
    \setlength{\columnseprule}{0.4pt}
    \setlength{\jot}{5pt}
    \setlength\parindent{0pt}
    \twocolumn
    \fcolorbox{white}{myGrey}{\begin{minipage}{20em}
        \centering \vspace{4mm} {\Large \textbf{Formula Sheet}} \\ 
        \theCourse\ (\theCourseCode) \\
        {\small \theExamDate} \vspace{3mm} 
    \end{minipage}}
}

% TITLES AND FORMATS
\pointname{} % No title for marks
\pointsdroppedatright % Marks at right margin
\framedsolutions % Add frame to solutions
\addto\captionsenglish{ % Format Figure labels
    \renewcommand\figurename{Figure}}
\addto\captionsenglish{ % Format Table labels
    \renewcommand\tablename{Table}}
\makeatletter
\def\fnum@figure{\figurename~\Alph{numberOfSections}\thefigure}
\def\fnum@table{\tablename~\Alph{numberOfSections}\thetable}
\makeatother
\renewcommand{\solutiontitle}{
   \noindent\textbf{Soluzione:}\par\noindent} % Rename solutions

% COVER PAGE - DO NOT EDIT
\newcommand{\theStudentName}{Mario/a Rossi}
\newcommand{\theStudentMatriculation}{0542185}
\newcommand{\theStudentIDPath}{support/unipacard.png}

\usepackage[nomessages]{fp}% http://ctan.org/pkg/fp

\newcommand*{\ExtractFirstChar}[2]{%
    \StrRemoveBraces{#2}[\FirstChar]%
    \StrChar{\FirstChar}{#1}[\FirstChar]%
    \FirstChar}
\newcommand{\theStudentAlpha}{
    \ExtractFirstChar{4}{\theStudentMatriculation}}
\newcommand{\theStudentBeta}{
    \ExtractFirstChar{5}{\theStudentMatriculation}}
\newcommand{\theStudentGamma}{
    \ExtractFirstChar{6}{\theStudentMatriculation}}

%% Prepare Email
\usepackage{etoolbox}
\makeatletter
\newcommand\myemail[3]{%                %\newcommand\tpj@compose@mailto[3]{%
\edef\@tempa{mailto:#1?subject=#2 }%
\edef\@tempb{\expandafter\html@spaces\@tempa\@empty}%
\href{\@tempb}{#3}}
\catcode\%=11
\def\html@spaces#1 #2{#1%20\ifx#2\@empty\else\expandafter\html@spaces\fi#2}
\catcode\%=14
\makeatother
\newcommand{\theExamSubjectCode}{\theExamAcronym\theExamDate}
\newcommand{\theEmailSubject}{[\theExamSubjectCode] \theStudentName, \theStudentMatriculation}







%%%\usepackage[pdftex]{hyperref}

\usepackage{etoolbox}
\makeatletter
\let\HyField@calcorder\ltx@empty
\Hy@AtBeginDocument{%
  \if@filesw
    \immediate\write\@mainaux{%
      \string\providecommand\string\HyField@AuxAddToCalculationOrder[1]{}%
    }%
  \fi
  \def\HyField@AuxAddToCalculationOrder#1{%
    \xdef\HyField@calcorder{%
      \ifx\HyField@calcorder\@empty
      \else
        \HyField@calcorder
        \space
      \fi
      #1 0 R%
    }%
  }%
}
\def\HyField@@AddToFields#1{%
  \HyField@AfterAuxOpen{%
    \if@filesw
      \write\@mainaux{%
        \string\HyField@AuxAddToFields{#1}%
      }%
      \ifx\Fld@calculate@code\@empty
      \else
        \write\@mainaux{%
          \string\HyField@AuxAddToCalculationOrder{#1}%
        }%
      \fi
    \fi
  }%
}%
\patchcmd{\@Form}{/Fields[\HyField@afields]}{/CO[\HyField@calcorder]/Fields[\HyField@afields]}{}{}
\makeatother




\def\theExamDate{\today}



\newcommand{\afcover}
{
\MakeResultCmd{7}%\theafQuestionCnt}
%\meaning\afResultCmd % uncomment to inspect

\begin{coverpages}
\begin{flushright}
{\Large
\begin{Form}
\TextField[readonly=true, value=0, charsize=16pt, backgroundcolor=verylightgray, calculate={\afResultCmd}]{{\bf Risultato:~}}
\end{Form}
\\
}
{\em (a cura del docente)}
\end{flushright}
\hrule
%
\begin{flushleft}
\Large
\begin{tabular}{r p{11cm}}
Esame di: & {\bf \theCourse}\\
Tema: & {\bf \theThemeTitle}\\
Consegna: & {\bf \theExamDate}\\[20pt]
%
Gruppo: & {\bf \theGroupID}\\
Studenti: & {\bf \theStudentOne} (corrispondente)\\
\ifdefined\theStudentTwo & {\bf \theStudentTwo}\fi\\
\ifdefined\theStudentThree & {\bf \theStudentThree}\fi\\
\ifdefined\theStudentFour & {\bf \theStudentFour}\fi\\
\ifdefined\theStudentFive & {\bf \theStudentFive}\fi\\
\ifdefined\theStudentSix & {\bf \theStudentSix}\fi\\
\ifdefined\theStudentSeven & {\bf \theStudentSeven}\fi\\
\ifdefined\theStudentEight & {\bf \theStudentEight}\fi\\
\end{tabular}
\end{flushleft}
%
\vfill
%
% \begin{flushright}
% \tcbox[enhanced,drop shadow,colback=white]{\myemail{\theExamEmail}{\theEmailSubject}{{\em Prepara email: \theExamEmail} } } (se non automaticamente compilato, indicare come soggetto \\dell'email: ``\theEmailSubject''; \\ includere manualmente il PDF della soluzione)
% \end{flushright}
\vfill 
% \begin{flushleft}
% {\large \bf Dichiarazione di buona condotta durante lo svolgimento dell'esame}\\ \vspace{3mm} {\em Io sottoscritto(a) {\bf \theStudentName}, studente(essa) dell'Università degli Studi di Palermo, con numero di matricola {\bf\theStudentMatriculation}, dichiaro di aver svolto la soluzione qui sotto riportata in completa autonomia e senza l'uso di materiale non consentito.}
% \end{flushleft}
\begin{flushright}
\small Powered by {\bf MIRPALab.it}~\href{https://www.mirpalab.it}{\faHome}~\href{https://www.facebook.com/mirpalab/}{\faFacebookSquare}~\href{https://www.instagram.com/mirpalab/}{\faInstagram}
\end{flushright}
\end{coverpages}
\newpage
}

\ifthenelse{\equal{\withSolutions}{Yes} \OR \equal{\withSolutions}{Y} \OR \equal{\withSolutions}{yes} \OR \equal{\withSolutions}{y} \OR \equal{\withSolutions}{YES}}{\printanswers}{}

%
\def\afpt{0}

% AF: dynamic macros
%\newcommand\MakeArray[2]{\ifnumless{#2}{1}{\def#1{}}{\def#1{~0=0} \ifnumgreater{#2}{1}{\foreach \n in {1,...,#2} {\xdef#1{#1,~\n=\n}}}{}}}
\newcommand\MakeArray[2]{\ifnumless{#2}{1}{\def#1{}}{\def#1{0} \ifnumgreater{#2}{1}{\foreach \n in {1,...,#2} {\xdef#1{#1,\n}}}{}}}
%
\usepackage{pgffor}
\newcommand{\Repeat}[2]{\foreach \n in {1,...,#1}{#2}}
\newcommand\MakeResultCmd[1]{
\newcounter{ct}\setcounter{ct}{2}
\def\afResultCmd{event.value = this.getField('answerA').value} 
\Repeat{#1}{\xdef\afResultCmd{\afResultCmd+this.getField('answer\Alph{ct}').value} \addtocounter{ct}{1}} \xdef\afResultCmd{\afResultCmd;}}
%
\newcounter{afQuestionCnt}\setcounter{afQuestionCnt}{0}
\newcommand\afnotess[1]
{
\stepcounter{afQuestionCnt}
\begin{flushright}
\begin{Form}
\noindent
\MakeArray\MyArray#1
\TextField[name=tanswer\Alph{afQuestionCnt}, multiline=true, bordercolor=verylightgray, backgroundcolor=verylightgray, width=0.95\linewidth, height=1in, charsize=12pt, color=myRed]{{\bf Correzione} {\em (a cura del docente)}:\\[1pt]}\\
%\ChoiceMenu[radio, name=answer\Alph{afQuestionCnt}, value=0]{{\em Punteggio:~}}{\MyArray}
\TextField[name=answer\Alph{afQuestionCnt}, bordercolor=verylightgray, backgroundcolor=verylightgray, width=0.05\linewidth, height=0.05\linewidth, charsize=12pt, color=blue, align=1]{{\em Punteggio:~}}\\
\end{Form}
\end{flushright}
}
\def\afnotes{\afnotess{\afpt}}

% \def\psimbA{\afnotes\def\afpt{0}\droppoints}
% \def\psimbB{\afnotes\def\afpt{0}\droppoint}
% \def\psimbC{\afnotes\def\afpt{0}\droppoints}
% \def\psimbD{\afnotes\def\afpt{0}\droppoints}
% \def\psimbE{\afnotes\def\afpt{0}\droppoints}
% \def\psimbF{\afnotes\def\afpt{0}\droppoints}

% \def\pnumA{\afnotes\def\afpt{0}\droppoints}
% \def\pnumB{\afnotes\def\afpt{0}\droppoint}
% \def\pnumC{\afnotes\def\afpt{0}\droppoints}
% \def\pnumD{\afnotes\def\afpt{0}\droppoints}
% \def\pnumE{\afnotes\def\afpt{0}\droppoints}
% \def\pnumF{\afnotes\def\afpt{0}\droppoints}

\def\psimbA{}
\def\psimbB{}
\def\psimbC{}
\def\psimbD{}
\def\psimbE{}
\def\psimbF{}

\def\pnumA{}
\def\pnumB{}
\def\pnumC{}
\def\pnumD{}
\def\pnumE{}
\def\pnumF{}

\def\pteorA{}

% \def\finalComments
% {
% \begin{Form}
% \noindent
% \TextField[name=tCommenti, multiline=true, bordercolor=verylightgray, backgroundcolor=verylightgray, width=0.95\linewidth, height=1in, charsize=12pt, color=myRed]{{\bf Commenti generali} {\em (a cura del docente)}:\\[1pt]}\\
% \end{Form}
% }

\def\finalComments{}




\def\afTheme{%
%
\section*{Tema {\em \theThemeTitle}}
\begin{multicols}{2} \noindent\theThemeDescription \thePicture \end{multicols}
\vfill\theModel\vfill 
\begin{multicols}{2} \theVariables\theParameters \end{multicols}
\theRequirements
\vfill \newpage
%
}



\def\afSubmission{%

\section*{Istruzioni per la consegna}
{\small 

\noindent Ogni tema progettuale si propone di far sperimentare allo studente, sia pure in modo semplificato, i diversi momenti di analisi e progettazione, sottesi alla soluzione di un problema di controllo; ciò avviene, per ovvie ragioni, attraverso uno studio teorico e solo simulato, come è tipico nei corsi di controlli automatici di base, ma rappresenta comunque un momento fondamentale che precede qualsiasi realizzazione completa e sperimentale. Le Tabelle~$1$ e~$2$ indicano il significato delle variabili e dei parametri del modello~(\ref{eq:model}), mentre le prima due righe della Tabella~$3$ specificano le ipotesi operative di equilibrio, mentre le restanti righe i requisiti di progetto.\\

La consegna del progetto avviene inviando un'email all'indirizzo di posta elettronica~\textit{\theExamEmail}, con oggetto ``[\theGroupID] Consegna'' e con allegato il PDF generato da questo documento. La consegna può essere fatta solo dal corrispondente del gruppo.\\

Dal momento della consegna non sarà più possibile modificare il seguente progetto overleaf, il quale dovrà includere almeno i seguenti contenuti:
\begin{itemize}
%
\item i file sorgenti \LaTeX{} del PDF (main.tex, theme.tex, 1*.tex, 2*tex, 3*.tex);
%
\item una cartella \textit{img} contenente le immagini e tutti i grafici inseriti nel testo, che devono essere riportati in formato vettoriale e contenere annotazioni adeguate e leggibili;
\item una cartella \textit{support} contenente (1) gli script (estensione .m) necessari per inizializzare e lanciare la simulazione, (2) i salvataggi delle sessioni di lavoro in \emph{sisotool}, e (3) i due diagrammi \emph{Simulink} (estensione .slx) dei sistemi approssimato lineare e non lineare, controllati (eventualmente con il doppio anello di retroazione), correttamente funzionanti e tali da riprodurre i risultati riportati nel documento descrittivo. 
\end{itemize}
\vfill \newpage

}

}


